% =====================================================================
% PREAMBULE COMMUN a cours_memoire.tex (parties I a IX) et a
% cours_memoire_questions.tex (partie X et suivantes). Extrait du fichier
% d'origine le 6 octobre 2026. Un seul endroit : toute retouche de style
% vaut pour les deux documents.
% =====================================================================
\documentclass[11pt,a4paper]{article}

\usepackage[utf8]{inputenc}
\usepackage[T1]{fontenc}
\usepackage{lmodern}
\usepackage[french]{babel}
\usepackage[margin=2.1cm]{geometry}
\usepackage{microtype}
\usepackage{amsmath,amssymb}
\usepackage{eurosym}
\usepackage{booktabs}
\usepackage{array}
\usepackage{tabularx}
\usepackage[table]{xcolor}
\usepackage{graphicx}
\usepackage{caption}
\usepackage{enumitem}
\usepackage[most]{tcolorbox}
\usepackage{tikz}
\usetikzlibrary{arrows.meta,positioning,calc,fit,backgrounds,decorations.pathreplacing,shapes.geometric}
\usepackage[hidelinks]{hyperref}

\definecolor{navy}{HTML}{1F3864}
\definecolor{bluemid}{HTML}{2E5496}
\definecolor{lightblue}{HTML}{D6E0F0}
\definecolor{accent}{HTML}{C0491F}
\definecolor{vert}{HTML}{3D8361}

\captionsetup{font=small,labelfont={bf,color=navy},labelsep=period}
\graphicspath{{../figures/}}

\newtcolorbox{phrasebox}{colback=vert!7,colframe=vert!70!black,boxrule=0.7pt,
  arc=2pt,left=7pt,right=7pt,top=5pt,bottom=5pt,before skip=6pt,after skip=8pt}
\newtcolorbox{cle}[1][]{colback=lightblue!40,colframe=navy,boxrule=0.6pt,
  arc=2pt,left=7pt,right=7pt,top=5pt,bottom=5pt,#1}

\providecommand{\Prob}{\mathbb{P}}
\providecommand{\E}{\mathbb{E}}
\providecommand{\R}{\mathbb{R}}

\newcommand{\lequoi}[1]{\smallskip\noindent\textcolor{navy}{\textbf{Le problème.}}\ #1\par}
\newcommand{\lidee}[1]{\smallskip\noindent\textcolor{navy}{\textbf{L'idée.}}\ #1\par}
\newcommand{\laformule}[1]{\smallskip\noindent\textcolor{navy}{\textbf{La formule.}}\ #1\par}
\newcommand{\cheznous}[1]{\smallskip\noindent\textcolor{navy}{\textbf{Chez nous.}}\ #1\par}
\newcommand{\lepiege}[1]{\smallskip\noindent\textcolor{accent}{\textbf{Le piège.}}\ #1\par}
\newcommand{\laphrase}[1]{\begin{phrasebox}\textbf{À dire à voix haute.} \emph{#1}\end{phrasebox}}

% Questions du jury : la question, puis la reponse a tenir.
\newtcolorbox{question}[1]{colback=accent!5,colframe=accent!75!black,boxrule=0.6pt,
  arc=2pt,left=7pt,right=7pt,top=5pt,bottom=5pt,before skip=7pt,after skip=4pt,
  fonttitle=\bfseries\color{white},title={#1}}
\newcommand{\rep}[1]{\smallskip\noindent\textcolor{navy}{\textbf{La réponse.}}\ #1\par}
\newcommand{\appui}[1]{\smallskip\noindent\small\textcolor{navy}{\textbf{Appui~:}} #1\par\normalsize}

% Partie X, chapitre par chapitre : les reperes de pages, le chapitre en une phrase, et la
% fiche \og par coeur\fg{} (tableau a deux colonnes dans un cadre, titre optionnel).
\newcommand{\reperes}[1]{\noindent{\small\color{navy!75}#1}\par\smallskip}
\newcommand{\enbref}[1]{\noindent\textcolor{navy}{\textbf{En une phrase.}}\ #1\par}
\newtcolorbox{parcoeur}[1][À savoir par c\oe ur]{colback=navy!3,colframe=navy!85!black,
  boxrule=0.6pt,arc=2pt,left=6pt,right=6pt,top=4pt,bottom=4pt,before skip=7pt,after skip=5pt,
  fonttitle=\bfseries\small\color{white},title={#1}}
\newenvironment{fichetab}{\small\renewcommand{\arraystretch}{1.1}%
  \tabularx{\linewidth}{@{}>{\raggedright\arraybackslash\bfseries}p{3.2cm}
  >{\raggedright\arraybackslash}X@{}}}{\endtabularx}

% Styles des schemas
\tikzset{
  brique/.style={draw=navy,fill=lightblue!45,rounded corners=2pt,inner sep=5pt,
                 align=center,font=\footnotesize,minimum height=9mm},
  pilier/.style={draw=navy,fill=white,circle,inner sep=1.5pt,minimum size=8mm,
                 font=\footnotesize\bfseries,text=navy},
  fleche/.style={-{Stealth[length=2mm]},draw=navy,line width=0.5pt},
  flechef/.style={-{Stealth[length=2.6mm]},draw=accent,line width=1.1pt},
  flechep/.style={-{Stealth[length=1.8mm]},draw=navy!45,line width=0.4pt},
  etiq/.style={font=\scriptsize\itshape,text=navy!75},
}

\setcounter{secnumdepth}{2}
\emergencystretch=1.2em   % noms propres longs (Kolmogorov-Smirnov, Pickands-Balkema)
